\section{Finite polylogarithmic reductions and mass antiderivatives}
\label{qspec:sec:finite}

\subsection{All nonpositive integer spectral points}
Put
\begin{equation}\label{qspec:eq:Lm}
L_m(a,b)=\left.\partial_s E(s;a,b)\right|_{s=-m},\qquad
z=\exp(-2\pi b+2\pi ia),\qquad m\ge0.
\end{equation}
The half-integer Bessel expression is the finite polynomial
\begin{equation}\label{qspec:eq:half-integer-K}
K_{m+1/2}(u)=\sqrt{\frac\pi{2u}}e^{-u}
\sum_{j=0}^m\frac{(m+j)!}{j!(m-j)!(2u)^j}.
\end{equation}
Its $m=0$ case and the three-term order recurrence prove it by induction;
it is also the classical half-integer formula \cite{halfK}.

\begin{theorem}[Finite Hurwitz--polylogarithm identities at every negative point]
\label{qspec:thm:finite}
For $0<a<1$, $b>0$, and every $m\ge0$,
\begin{align}
L_m(a,b)={}&2\Rea\sum_{j=0}^m\binom mj(-2ib)^{m-j}
\bigl\{\zeta'(-m-j,a+ib)+\zeta'(-m-j,1-a+ib)\bigr\}\label{qspec:eq:L-Hurwitz}\\
={}&-\frac{2^{2m+1}(m!)^2\pi}{(2m+1)!}b^{2m+1}\notag\\
&+2(-1)^mm!\sum_{j=0}^m
\frac{(m+j)!b^{m-j}}{j!(m-j)!4^j\pi^{m+j}}
\Rea\Li_{m+j+1}(z).
\label{qspec:eq:L-polylog}
\end{align}
In particular, the equality of the two right sides is a finite special-function
identity; $L_m$ is not needed as an additional unevaluated coordinate.
\end{theorem}
\begin{proof}
The reflected decomposition of $E$ is
$E(s;a,b)=Z_a(s;b)+Z_{1-a}(s;b)$. Apply
Corollary~\ref{qspec:cor:negative-hurwitz} to the two terms to obtain
\eqref{qspec:eq:L-Hurwitz}.

For the independent evaluation, write the Poisson expression as
$E(s)=H(s)/\Gamma(s)$, where $H$ is holomorphic at $s=-m$. The zero
of the reciprocal Gamma function is simple and
\[
(1/\Gamma)'(-m)=(-1)^mm!.
\]
Consequently $E'(-m)=(-1)^mm!H(-m)$: derivatives of the Bessel order
and of the other factors in $H$ do not appear in this first jet. The
zero Fourier mode becomes
\[
(-1)^mm!\sqrt\pi\Gamma(-m-\tfrac12)b^{2m+1}
=-\frac{2^{2m+1}(m!)^2\pi}{(2m+1)!}b^{2m+1}.
\]
For the nonzero modes use $K_{-m-1/2}=K_{m+1/2}$ and
\eqref{qspec:eq:half-integer-K}. The square-root prefactor reduces to
$1/(2\sqrt{bn})$. The remaining $n$ powers are
$n^{-m-j-1}$, so their cosine sums are exactly
$\Rea\Li_{m+j+1}(z)$. Collecting the powers of $b$, $4$, and $\pi$
gives \eqref{qspec:eq:L-polylog}.
\end{proof}

The first four values, writing $\Rea\Li_j=\Rea\Li_j(z)$, are
\begin{align}
L_0={}&-\log\bigl(2(\cosh(2\pi b)-\cos(2\pi a))\bigr),\label{qspec:eq:L0}\\
L_1={}&-\frac{4\pi b^3}{3}-\frac{2b}{\pi}\Rea\Li_2
                       -\frac1{\pi^2}\Rea\Li_3,\label{qspec:eq:L1}\\
L_2={}&-\frac{16\pi b^5}{15}+\frac{4b^2}{\pi^2}\Rea\Li_3
              +\frac{6b}{\pi^3}\Rea\Li_4+\frac3{\pi^4}\Rea\Li_5,\label{qspec:eq:L2}\\
L_3={}&-\frac{32\pi b^7}{35}-\frac{12b^3}{\pi^3}\Rea\Li_4
              -\frac{36b^2}{\pi^4}\Rea\Li_5\notag\\
     &\hspace{19mm}-\frac{45b}{\pi^5}\Rea\Li_6
              -\frac{45}{2\pi^6}\Rea\Li_7.
\label{qspec:eq:L3}
\end{align}
For $m=0$, $\Li_1(z)=-\log(1-z)$ reduces the finite formula to
\eqref{qspec:eq:L0}. For $m=1$, the Hurwitz side is exactly
\eqref{qspec:eq:opening-example}.

\subsection{An exact Appell-type ladder in the squared mass}
Here the term ``Appell-type'' means the displayed derivative relation,
not an assertion that these functions are polynomials.

\begin{theorem}[Mass differentiation and every anchored mass primitive]
\label{qspec:thm:mass-ladder}
With $u=b^2$,
\begin{equation}\label{qspec:eq:mass-derivative}
\partial_u L_m(a,\sqrt u)=mL_{m-1}(a,\sqrt u)\quad(m\ge1),
\qquad
\partial_u L_0(a,\sqrt u)=-E(1;a,\sqrt u).
\end{equation}
The continuous zero-mass values are
\begin{align}
L_0(a,0)&=-\log(4\sin^2\pi a),\label{qspec:eq:massless0}\\
L_m(a,0)&=\frac{2(-1)^m(2m)!}{(2\pi)^{2m}}
\Rea\Li_{2m+1}(e^{2\pi ia})\quad(m\ge1).
\label{qspec:eq:masslessm}
\end{align}
For every $m\ge1$,
\begin{align}
L_m(a,b)={}&\sum_{j=0}^{m-1}\binom mj b^{2j}L_{m-j}(a,0)\notag\\
&+m\int_0^{b^2}(b^2-u)^{m-1}L_0(a,\sqrt u)\dd u.
\label{qspec:eq:mass-primitives}
\end{align}
Thus the finite polylogarithm formulas evaluate every integral on the
right side exactly.
\end{theorem}
\begin{proof}
Initially in an absolute-convergence half-plane,
$\partial_u E(s)=-sE(s+1)$. Continue meromorphically and differentiate
in $s$:
\[
\partial_u E'(s)=-E(s+1)-sE'(s+1).
\]
At $s=-m$, $m\ge1$, the first term is zero because
$E(1-m)=0$. This proves the first relation. At $s=0$ it gives the
second relation.

For fixed $a\in(0,1)$ there is a nonzero gap between the lattice and
zero. The binomial continuation, after finitely many Hurwitz
subtractions if needed, is analytic in $u$ near zero. Alternatively,
take the limit in the explicit finite formula: for $m\ge1$ only its
$j=m$ term survives, while the absolutely convergent polylogarithms
have continuous radial limits. This proves
\eqref{qspec:eq:masslessm}; \eqref{qspec:eq:massless0} follows directly
from \eqref{qspec:eq:L0}. Repeated integration of
\eqref{qspec:eq:mass-derivative}, with these initial values, gives
\eqref{qspec:eq:mass-primitives}. The coefficient of its integral is
$m!/(m-1)!=m$.
\end{proof}

For example, a non-elementary logarithmic integral is evaluated by
\begin{align}
&-\int_0^{b^2}\log\bigl(2(\cosh(2\pi\sqrt u)-\cos(2\pi a))\bigr)\dd u\notag\\
&\qquad=-\frac{4\pi b^3}{3}-\frac{2b}{\pi}\Rea\Li_2(z)
-\frac1{\pi^2}\Rea\Li_3(z)
+\frac1{\pi^2}\Rea\Li_3(e^{2\pi ia}).
\label{qspec:eq:log-mass-integral}
\end{align}
The higher-weight analogues follow by inserting $m=2,3,\ldots$ in
\eqref{qspec:eq:mass-primitives}. Their constants are fixed by the
zero-mass boundary values, not left as undetermined integration terms.

The same ladder provides an independent algebraic check of the finite
coefficients. Since $\partial_b\Li_j(z)=-2\pi\Li_{j-1}(z)$,
substitution of \eqref{qspec:eq:L-polylog} into
$\partial_bL_m=2bmL_{m-1}$ leaves a finite identity of rational
coefficients and powers of $b,\pi$. The delivered exact suite checks
these cancellations through $m=10$; the theorem itself holds for all
$m$ by the analytic proof.

\subsection{A differential equation in the shift and mass}
The same finite family is closed under a second differential relation.
\begin{proposition}\label{qspec:prop:PDE}
For every $m\ge0$,
\begin{equation}\label{qspec:eq:PDE}
\left(\partial_a^2+\partial_b^2-\frac{2m}{b}\partial_b\right)L_m(a,b)=0.
\end{equation}
Equivalently, with $u=b^2$,
\[
\bigl(\partial_a^2+4u\partial_u^2+(2-4m)\partial_u\bigr)
L_m(a,\sqrt u)=0.
\]
For $m\ge2$ this becomes the finite adjacent-weight identity
\begin{equation}\label{qspec:eq:PDE-adjacent}
\partial_a^2L_m=2m(2m-1)L_{m-1}-4m(m-1)b^2L_{m-2}.
\end{equation}
\end{proposition}
\begin{proof}
Direct differentiation of one radial summand gives
\[
(\partial_a^2+\partial_b^2)E(s)=4s^2E(s+1),
\qquad \partial_bE(s)=-2sbE(s+1).
\]
Hence $(\partial_a^2+\partial_b^2+(2s/b)\partial_b)E(s)=0$
by continuation from absolute convergence. Differentiate in $s$ at
$s=-m$. The extra term $(2/b)\partial_bE(-m)$ is zero, because
$E(-m)$ vanishes identically as a function of the parameters. This
proves \eqref{qspec:eq:PDE}. The mass ladder yields
\eqref{qspec:eq:PDE-adjacent} for $m\ge2$.
\end{proof}
For $m=0$ the logarithm in \eqref{qspec:eq:L0} is harmonic on the open
strip $b>0$. The higher weights satisfy its explicit weighted version,
rather than an unrelated differential equation at each polylogarithm order.

\subsection{Products of several quadratic factors}
The parity mechanism extends beyond a single conjugate pair. The following
finite reduction uses no multivariable special functions at the first
spectral derivative.

\begin{theorem}[Several masses with an arbitrary even polynomial weight]
\label{qspec:thm:several-masses}
Let $P$ be a polynomial of degree at most $D$, let $r\ge1$, and let
$b_1,\ldots,b_r>0$. Define, initially in its absolute-convergence
half-plane,
\begin{equation}\label{qspec:eq:several-masses}
\mathcal Z_{P,\mathbf b}(s;a)=
\sum_{n\in\Z}P((n+a)^2)
\prod_{j=1}^r\bigl((n+a)^2+b_j^2\bigr)^{-s}.
\end{equation}
Its continuation is holomorphic at zero, and
\begin{equation}\label{qspec:eq:several-masses-finite}
\mathcal Z_{P,\mathbf b}'(0;a)
=\sum_{j=1}^r\sum_{\ell=0}^D
\frac{P^{(\ell)}(-b_j^2)}{\ell!}L_\ell(a,b_j).
\end{equation}
Combining this formula with \eqref{qspec:eq:L-polylog} gives a finite
polylogarithm reduction for every such spectrum, including repeated masses.
\end{theorem}
\begin{proof}
First work on a positive tail $x=n+a$ with all $b_j<x$. After factoring
out $x^{-2rs}$, the remaining factor is
\[
\exp\left(-s\sum_{j=1}^r\log(1+b_j^2/x^2)\right).
\]
It has only even powers of $x^{-1}$. Multiplication by $P(x^2)$ still
leaves only even powers of $x$. A possible Hurwitz pole at $s=0$
would require the resulting power to be $x^{-1}$, which is impossible.
Thus there are no resonant pole terms at zero in this common-scale
expansion.

Compare the product spectrum with $r^{-1}$ times the sum of the
individual weighted spectra evaluated at $rs$. Their constant and
linear terms in $s$ agree identically before summation; the remaining
coefficients are $O(s^2)$. Since there is no pole at zero, their first
derivatives agree after continuation as well. Remove a finite head
to justify the expansion for arbitrary masses. The ordinary head
logarithms add, so it contributes no discrepancy. The same argument
applies to the negative lattice tail, where evenness changes its
weight to $P(y_n^2)$ without a sign.

Finally the polynomial Taylor identity
\[
P(x^2)=\sum_{\ell=0}^D\frac{P^{(\ell)}(-b_j^2)}{\ell!}
 (x^2+b_j^2)^\ell
\]
expresses the $j$th individual weighted derivative as the inner finite
sum in \eqref{qspec:eq:several-masses-finite}. This proves the formula.
\end{proof}

For instance, $P(v)=v+c$ and two masses give
\[
\mathcal Z_{P,(b_1,b_2)}'(0;a)
=\sum_{j=1}^2\bigl(L_1(a,b_j)+(c-b_j^2)L_0(a,b_j)\bigr).
\]
Each summand has only elementary logarithms, $\Li_2$, and $\Li_3$.
The theorem is restricted to the first spectral derivative. Products
of logarithms in higher derivatives are not removed by this argument.

\subsection{Positive-integer resolvents: polygamma and rational polylogs}
At positive integers the same half-integer Bessel polynomial produces
polylogarithms of nonpositive order.

\begin{proposition}\label{qspec:prop:positive}
For $N\ge1$,
\begin{align}
E(N;a,b)={}&\frac{\sqrt\pi\Gamma(N-1/2)}{\Gamma(N)}b^{1-2N}\notag\\
&+\frac{2\pi^N b^{-N}}{\Gamma(N)}
\sum_{j=0}^{N-1}\frac{(N-1+j)!}
 {j!(N-1-j)!(4\pi b)^j}\Rea\Li_{j+1-N}(z).
\label{qspec:eq:positive-polylog}
\end{align}
An independent finite polygamma expression is
\begin{equation}\label{qspec:eq:positive-polygamma}
E(N;a,b)=2\Rea\sum_{r=1}^N A_{N,r}(b)
\bigl(\Xi_r(a+ib)+\Xi_r(1-a+ib)\bigr),
\end{equation}
where
\begin{align}
A_{N,r}(b)&=(-1)^{N-r}\binom{2N-r-1}{N-r}(-2ib)^{-(2N-r)},\label{qspec:eq:Ar}\\
\Xi_1(c)&=-\psi(c),\qquad
\Xi_r(c)=\frac{(-1)^r}{(r-1)!}\psi^{(r-1)}(c)\quad(r\ge2).
\label{qspec:eq:Xi}
\end{align}
\end{proposition}
\begin{proof}
Insert \eqref{qspec:eq:half-integer-K}, with $m=N-1$, into
\eqref{qspec:eq:poisson}. The remaining powers of $n$ give
\eqref{qspec:eq:positive-polylog}.

For the other formula, expand the rational function at $x=-ib$:
\[
\frac1{(x^2+b^2)^N}
=\sum_{r=1}^N\left\{\frac{A_{N,r}(b)}{(x+ib)^r}
+\frac{\overline{A_{N,r}(b)}}{(x-ib)^r}\right\}.
\]
The coefficient follows by expanding $(y-2ib)^{-N}$ at $y=x+ib=0$.
For $r\ge2$ summation gives $\zeta(r,c)=\Xi_r(c)$. The $r=1$
terms must be paired before summation: their coefficients add to zero,
and a common finite cutoff gives finite part $-\psi(c)$. No separate
value for a divergent harmonic series is assigned. Sum over $x_n$
and $y_n$ to obtain \eqref{qspec:eq:positive-polygamma}.
\end{proof}
For $N=1$ these formulas reduce to the classical resolvent
\begin{equation}\label{qspec:eq:resolvent}
E(1;a,b)=\frac\pi b
\frac{\sinh(2\pi b)}{\cosh(2\pi b)-\cos(2\pi a)}.
\end{equation}
For larger $N$, all polylogs in
\eqref{qspec:eq:positive-polylog} are rational functions of $z$, obtained
by repeatedly applying $z\partial_z$ to $z/(1-z)$.

\subsection{Gamma products and every mass primitive of their logarithm}
Define the real quantity
\begin{equation}\label{qspec:eq:Rgamma}
R(a,b)=\log\frac{\Gamma(a)^2\Gamma(1-a)^2}
 {\Gamma(a+ib)\Gamma(a-ib)\Gamma(1-a+ib)\Gamma(1-a-ib)}.
\end{equation}
All conjugate products in the denominator are positive real numbers.

\begin{proposition}[A Gamma-product dilogarithm integral]
\label{qspec:prop:gamma-integrals}
The classical product and reflection identities give
\begin{equation}\label{qspec:eq:Rhyperbolic}
R(a,b)=\log\frac{\cosh(2\pi b)-\cos(2\pi a)}{1-\cos(2\pi a)}.
\end{equation}
In particular, with $z_0=e^{2\pi ia}$,
\begin{equation}\label{qspec:eq:gamma-Li2}
\int_0^b R(a,t)\dd t
=\pi b^2-b\log(4\sin^2\pi a)
+\frac1\pi\Rea\bigl(\Li_2(z)-\Li_2(z_0)\bigr).
\end{equation}
More generally, for every $r\ge1$,
\begin{align}
&\frac1{(r-1)!}\int_0^b(b-t)^{r-1}R(a,t)\dd t\notag\\
&\quad=\frac{2\pi b^{r+1}}{(r+1)!}
-\frac{b^r}{r!}\log(4\sin^2\pi a)\notag\\
&\qquad\ -2\Rea\left\{\frac1{(-2\pi)^r}
\left(\Li_{r+1}(z)-\sum_{j=0}^{r-1}\frac{(-2\pi b)^j}{j!}
\Li_{r+1-j}(z_0)\right)\right\}.
\label{qspec:eq:gamma-all-primitives}
\end{align}
\end{proposition}
\begin{proof}
The Gamma product, with its conjugate linear terms cancelling, gives
\[
\sum_{n\ge0}\log\left(1+\frac{b^2}{(n+a)^2}\right)
=2\log\Gamma(a)-\log\Gamma(a+ib)-\log\Gamma(a-ib).
\]
Add the formula at $1-a$ and use Gamma reflection \cite{gamma}; equivalently
compare \eqref{qspec:eq:L0} at $b$ and at zero. The result is
\eqref{qspec:eq:Rhyperbolic}. With the logarithmic branches already fixed,
it can be written
\[
R(a,b)=2\pi b-\log(4\sin^2\pi a)-2\Rea\Li_1(e^{-2\pi b}z_0).
\]
One integration using $\partial_b\Li_2(z)=-2\pi\Li_1(z)$ proves
\eqref{qspec:eq:gamma-Li2}. For repeated integration, an $r$th primitive
of $\Li_1(z)$ is $(-2\pi)^{-r}\Li_{r+1}(z)$. Subtract its Taylor
polynomial through order $r-1$ at zero to enforce all anchored initial
conditions. This gives \eqref{qspec:eq:gamma-all-primitives}. Unit-circle
values appearing in that Taylor polynomial have order at least two and
are absolutely convergent.
\end{proof}
